\documentclass[12pt,a4paper]{article}
%per compilare il template usare questo sotto e commentare l'altro:
%\usepackage{../stile}
\usepackage{../../../template/stile}

%Titolo documento
\newcommand{\titoloDocumento}{Studio di Fattibilià}

%Prima data di creazione del documento
\newcommand{\dataCreazione}{30 Novembre 2015}

%Inserite la versione attuale del documento
\newcommand{\versione}{1.2.0}

%Stato in cui si trova il documento: Formale solo all'atto di consegna
\newcommand{\stato}{Formale}

%Uso del documento
\newcommand{\uso}{Interno}

\rhead{\titoloDocumento}
\lfoot{Versione: \versione}
\title{\titoloDocumento}

\begin{document}
\begin{titlepage}
\begin{center}
\today \\
\vspace{1cm}
\begin{Huge}
\textbf{\nomeGruppo} \\
\end{Huge}
\textbf{\prjL} \\
\vspace{1cm}
\includegraphics[scale=0.5]{\logoGrande}
\vspace{1cm}

\HRule \\[0.4cm]
\begin{Huge}
{\huge \bfseries \titoloDocumento}\\[0.4cm]
\end{Huge}
\HRule \\[1cm]
\vfill

\begin{table}[h]
\begin{center}
\begin{tabular}{r | l}
\multicolumn{2}{c}{\textbf{Informazioni sul documento}}\\
\midrule
\textbf{Nome Documento} & \titoloDocumento \\
\textbf{Versione} & \versione \\
\textbf{Stato} & \emph{\stato} \\
\textbf{Uso} & \emph{\uso} \\
\textbf{Data Creazione} & \dataCreazione \\
\textbf{Data Ultima Modifica} & \today \\
\textbf{Redazione} & \IB\\
\  & \TP\\
\  &  \AVE\\
\textbf{Verifica} &  \NDC\\
\ & \AVI \\
\textbf{Approvazione} &  \WS\\
\textbf{Lista Distribuzione} & \nomeGruppo \\
\  & \Vardanega \\
\  & \Cardin \\
\  & Il proponente \Zucchetti \\

\end{tabular}
\end{center}
\end{table}

\end{center}
\end{titlepage}
\newpage

\Large{\textbf{Registro delle modifiche}}\\
\normalsize

\begin{table}[h]
\begin{center}

\begin{tabular}{p{0.12\textwidth} p{0.25\textwidth} p{0.18\textwidth} p{0.4\textwidth}}
\toprule
\textbf{Versione} & \textbf{Autore} & \textbf{Data} & \textbf{Descrizione}\\
\midrule
\midrule
v1.2.0  & \WS & 2015-12-21 & Correzioni minori ed approvazione\\
\midrule
v1.1.1  & \AVI & 2015-12-21 & Revisione: correzioni sintattiche\\
\midrule
v1.1.0  & \NDC & 2015-12-20 & Prima revisione: correzioni lessicali, grammaticali e di stile\\
\midrule
v1.0.7  & \TP & 2015-12-18 & Aggiunta introduzione\\
\midrule
v1.0.6  & \TP & 2015-12-18 & Normalizzate sezioni capitolati\\
\midrule
v1.0.5 & \AVE & 2015-12-08 & Aggiunto capitolato 2 e 6\\
\midrule
v1.0.4 & \IB & 2015-12-08 & Aggiunto capitolato 4\\
\midrule
v1.0.3 & \IB & 2015-12-08 & Aggiunto capitolato 3\\
\midrule
v1.0.2 & \TP & 2015-12-07 & Aggiunto capitolato 5\\
\midrule
v1.0.1 & \TP & 2015-12-06 & Aggiunto capitolato 1\\
\midrule
v1.0.0 & \NDC & 2015-11-30  & Creazione documento\\
\bottomrule
\end{tabular}
\caption{Versionamento del documento}
\label{tabVers1}
\end{center}
\end{table}
\newpage

\tableofcontents
\listoftables
\newpage

\section{Introduzione}

\subsection{Scopo}
In questo documento vengono analizzati i requisiti, aspetti positivi e fattori di rischio di ogni capitolato proposto per evidenziare i punti di forza e gli aspetti negativi che hanno portato a scegliere o scartare un capitolato.

\subsection{Descrizione}
Per il capitolato scelto dal gruppo ovvero \gloss{Quizzpedia} viene presentata una descrizione e una valutazione approfondita sugli aspetti positivi, negativi e i fattori di rischio che sono stati valutati e hanno determinato la scelta di questo capitolato. Gli altri capitolati vengono presentati con una breve descrizione. Per ognuno di essi vengono inoltre elencati aspetti positivi tenuti in considerazione e fattori di rischio che hanno spinto il gruppo ad escludere i capitolati in esame.

\subsection{Riferimenti}
\subsubsection{Informativi}

\begin{itemize}
\item \href{http://www.math.unipd.it/~tullio/IS-1/2015/Progetto/C1.pdf}{Capitolato C1}
\item \href{http://www.math.unipd.it/~tullio/IS-1/2015/Progetto/C2.pdf}{Capitolato C2}
\item \href{http://www.math.unipd.it/~tullio/IS-1/2015/Progetto/C3.pdf}{Capitolato C3}
\item \href{http://www.math.unipd.it/~tullio/IS-1/2015/Progetto/C4.pdf}{Capitolato C4}
\item \href{http://www.math.unipd.it/~tullio/IS-1/2015/Progetto/C5.pdf}{Capitolato C5}
\item \href{http://www.math.unipd.it/~tullio/IS-1/2015/Progetto/C6.pdf}{Capitolato C6}
\end{itemize}


\newpage
\section{Capitolato C5 - Quizzpedia}
\subsection{Descrizione}

Il capitolato propone lo sviluppo di un sito web per la creazione e somministrazione di questionari.

\subsubsection{Dominio applicativo}

Le domande verranno fornite attraverso uno specifico \textit{Quiz Markup Language} (QML), la cui creazione è uno degli oggetti principali del capitolato. L'interfaccia amministrativa dovrà supportare dispositivi \textit{desktop}, mentre per i candidati viene richiesto il supporto di dispositivi \textit{desktop} e \textit{mobile}. Si dovranno considerare tre tipi di profili:

\begin{itemize}
\item \textbf{Amministratore:} crea profili di tipo docente e ha il completo controllo e accesso alla piattaforma;
\item \textbf{Docente:} crea questionari, classi di studenti e analizza i risultati delle risposte;
\item \textbf{Studente:} si iscrive a classi e completa i questionari.
\end{itemize}

\subsubsection{Dominio tecnologico}

Per quanto riguarda il \textit{server} è richiesto l'utilizzo di una delle seguenti opzioni:

\begin{itemize}
\item Tomcat (Java);
\item Node.js (JavaScript).
\end{itemize}
Per l'archiviazione sono proposte le seguenti tecnologie:
\begin{itemize}
\item PostgreSQL;
\item MongoDB.
\end{itemize}
Per il \textit{client} viene richiesto l'utilizzo di HTML5, CSS3 e il supporto \textit{desktop} per l'intera applicazione e \textit{mobile} solo per la parte destinata allo studente.

Per la creazione del linguaggio QML il gruppo ha considerato l'utilizzo di XML, di Markdown o di altri linguaggi di markup estendibili.

\subsection{Valutazione}
\subsubsection{Aspetti positivi}

Sono stati evidenziati i seguenti aspetti positivi che hanno contribuito alla scelta di questo capitolato:

\begin{itemize}
\item semplicità e chiarezza degli obiettivi;
\item la gestione dell'interfaccia \textit{responsive} può essere risolto attraverso i numerosi \textit{framework} (tra cui \href{http://getbootstrap.com}{Bootstrap} e \href{http://foundation.zurb.com}{Foundation});
\item PostgreSQL e il sistema relazionale è propedeutico per questo corso quindi è conosciuto dall'intero gruppo;
\item HTML, CSS e Javascript sono trattate nel corso di Tecnologie Web che tutti i membri del gruppo stanno o hanno
seguito;
\item il cuore applicativo è semplice da implementare ed è possibile utilizzare un approccio incrementale
per lo sviluppo delle varie funzioni, sia richieste che opzionali.
\end{itemize}

\subsubsection{Potenziali criticità}

Sono state inoltre rilevate possibili punti critici nello sviluppo applicativo e/o nelle tecnologie utilizzate
che dovranno essere tenuti sotto controllo:

\begin{itemize}
\item nel caso di utilizzo di MongoDB, la tecnologia il paradigma (NoSQL) non sono affatto conosciuti dai membri del gruppo;
\item utilizzando Javascript bisogna tenere a mente che la semplicità e la dinamicità del linguaggio sono armi a doppio taglio;
\item le librerie Tomcat e Node.js sono sconosciute a tutti i membri del gruppo;
\end{itemize}
Per lo sviluppo Javascript, vengono consigliati strumenti per un maggiore controllo della correttezza del codice (come \href{http://flowtype.org}{Flow}) oppure altri linguaggi tipati che compilati producono Javascript (come \href{http://coffeescript.org}{CoffeeScript}).

\subsection{Conclusione}

Nonostante la scelta delle tecnologie da utilizzare verrà affrontata più avanti, le possibili scelte fornite dal capitolato sono tecnologie molto utilizzate nelle aziende per cui il loro utilizzo ci permetterebbe di apprendere nuove conoscenze e competenze spendibili nel mondo lavorativo.

Il dominio applicativo risulta chiaro e più in linea con il tipo di progetti a cui siamo stati sottoposti da altri corsi. Inoltre il capitolato è abbastanza flessibile da permettere di limitare alcuni rischi scegliendo strumenti più sicuri e familiari e di applicare direttamente conoscenze apprese dal corso di ingegneria del software quali un approccio incrementale per le varie funzionalità dell'applicazione.

Per questi motivi crediamo che gli obiettivi del capitolato e le tecnologie utilizzate siano in linea con gli obbiettivi formativi e le competenze che vogliamo acquisire.

\newpage
\section{Altri Capitolati}

\subsection{Capitolato C1 Actorbase}
\subsubsection{Descrizione}

Il capitolato richiede l'implementazione di un \textit{database} NoSQL\ di tipo \textit{key-value} basato
sull'architettura ad attori. E' inoltre richiesta la definizione di un \textit{domain-specific-language} (DSL)
per poter interagire con il \textit{database} da linea da comando.

\subsubsection{Aspetti positivi}
\begin{itemize}
\item Scala risulta un linguaggio interessante e conciso;
\item l'architettura proposta non è di difficile comprensione;
\item il modello \textit{actor-model} è presentato anche da Programmazione Concorrente e Distribuita di quest'anno, corso che è seguito dalla maggior parte del gruppo;
\item possibilità di scegliere Java al posto di Scala, considerando che tutto il gruppo ha esperienza con il primo;
\item il proponente interno potrebbe avere una maggiore disponibilità.
\end{itemize}

\subsubsection{Fattori di rischio}
\begin{itemize}
\item la libreria Akka è per lo più sconosciuta ai membri del gruppo;
\item nel caso si scelga Java, l'interazione con la libreria Akka risulterebbe di gran lunga più verbosa
in quanto essa è stata pensata per Scala;
\item poca conoscenza delle caratteristiche dei sistemi distribuiti;
\item lo sviluppo di un DSL appropriato potrebbe richiedere conoscenze relative al \textit{parsing} e 
analisi lessicale;
\item se Scala è scelto come linguaggio implementativo, sarà necessario imparare il linguaggio e l'ecosistema
di strumenti e librerie a sua disposizione.
\end{itemize}

\newpage

\subsection{Capitolato C2 CLIPS}

\subsubsection{Descrizione}

Per questo capitolato è richiesta l'implementazione di un'applicazione per \textit{smartphone}. L'applicazione deve permettere di gestire la micro-localizzazione dove è presente la tecnologia BLE Beacon, e veicolare all' utente finale informazioni contestuali a questi luoghi. 

\subsubsection{Aspetti positivi}
\begin{itemize}
 \item il gruppo ha trovato questo capitolato indubbiamente interessante perché posto nell'ambito della ricerca;
 \item permetterebbe lo studio di tecnologie recenti;
 \item la possibilità di proporre e sviluppare una propria idea originale.
\end{itemize}

\subsubsection{Fattori di rischio}
\begin{itemize}
\item la teconologia dei Beacons BLE e le sue applicazioni sono ancora immature;
\item abbiamo ritenuto che un costo medio di 25 euro per Beacon fosse eccessivo mettendolo in relazione con le funzionalità offerte;
\item l'obbligo da parte del sistema operativo Android o IOS di tenere attivo il sistema di localizzazione dello \textit{smartphone} per poter localizzare i Beacons risulta in un consumo notevole della batteria;
\item l'impossibilità di interagire in alcun modo con i Beacons senza un applicazione dedicata.
\item difficoltà nel proporre un idea non banale e funzionale considerando l'imprecisione e il tempo di latenza nel rilevamento dei Beacons da parte si uno \textit{smartphone}.
\end{itemize}

\newpage
\subsection{Capitolato C3 UMAP}

\subsubsection{Descrizione}

Questo capitolato richiede l'implementazione di un sistema per la  raccolta e la elaborazione di dati eterogenei, dando inoltre la possibilità all'utente finale di leggere i dati e l'analisi fatta sugli stessi da un \textit{engine} predittivo.

Il capitolato richiede lo sviluppo di:
\begin{itemize}
\item frontend;
\item business Logic API; 
\item \textit{engine} predittivo;
\item broker MQTT. 
\end{itemize}

\subsubsection{Aspetti positivi}
\begin{itemize}
\item le tecnologie necessarie per svolgere il capitolato sono svariate (Java, MongoDB, AWS , HTML e molte altre) il che permetterebbe al gruppo di venire in contatto con strumenti nuovi ed estremamente utili in ambito lavorativo.
\end{itemize}

\subsubsection{Fattori di rischio}

\begin{itemize}
\item lo studio e l'apprendimento delle tecnologie necessarie sicuramente risulterebbe in uno sforzo notevole da parte di ogni componente del gruppo;

\item un altra difficoltà risiede nella realizzazione del \textit{engine} predittivo dato che al gruppo mancano totalmente le idee e le conoscenze per concepire l'elaborazione di dati eterogenei;
\item il sistema da realizzare sarebbe incredibilmente amplio.
\end{itemize}

\newpage
\subsection{Capitolato C4 MaaS}
\subsubsection{Descrizione}
Maas è basato su un prodotto di nome Maap, che è un database NoSql per il trattamento di dati \textit{business} attraverso un interfaccia grafica permettendone l'utilizzo a persone che non sono esperte di \textit{database}. Il committente richiede che questo prodotto diventi un servizio gestito attraverso un sito web che dia la possibilità di eliminare l'onere all'utente di eseguire la preparazione e l'installazione di Maap.

Sono stati rilevati i seguenti requisiti:

\begin{itemize}

\item le tecnologie richieste includono:

\begin{itemize}
\item Node.js per il backend, MaaS deve essere in Long Term Support [LTS] versione Argon;
\item mongoDB versione maggiore di 3.x come database per l'applicazione; 
\end{itemize}

\item il Sistema deve essere costruito usando un \textit{framework} di alto livello per Node.js (noi fortemente consigliamo di usare IBM loopback);

\item il sistema deve essere schierabile con Heroku;
 
\item il codice sorgente deve essere pubblicato e versionato usando GitHub oppure Bitbucket. 

\end{itemize}

\subsubsection{Aspetti positivi}
\begin{itemize}
\item potenzialmente la mole di lavoro e la complessità dello stesso potrebbero essere esigue dato che il capitolato si basa già su un prodotto esistente che deve essere prima compreso e studiato in tutte le sue parti;

\item alcune delle tecnologie a supporto del capitolato potrebbero essere , anche se in modo parziale, già note al gruppo (HTML, CSS, DBMS, Java).
\end{itemize}

\subsubsection{Fattori di rischio}
\begin{itemize}
\item l'argomento risulta essere poco stimolante, dato che la parte più interessante sarebbe stata la creazione di Maap che è già stato implementato;

\item essendo il committente una \textit{startup} con risorse di tempo, personale e finanze limitate è possibile che sia difficile aiutare e seguire in gruppo di lavoro in caso di problemi o dubbi inerenti al progetto.
\end{itemize}

\newpage
\subsection{Capitolato C6 MIVOQ}
\subsubsection{Descrizione}

Il Capitolato richiede che il gruppo di lavoro proponga un idea e realizzi applicazione mobile che sfrutti le funzionalità di \textit{text-to-speech} fornite attraverso delle API via HTTP.

\subsubsection{Aspetti positivi}
\begin{itemize}
 \item le tecnologie necessarie per la creazione dell'applicazione mobile sono: su Android Java e la libreria android-SDK, su IOS Objective-C, su WindowsPhone .NET. Tutti questi linguaggi e \textit{framework} sono largamente documentati e usati, elemento che faciliterebbe sicuramente il lavoro del gruppo;
 \item il capitolato può essere molto stimolante grazie al fatto che il proponete da la possibilità al gruppo di sviluppare un'idea originale per l'applicazione. 
\end{itemize}
\subsubsection{Fattori di rischio}
\begin{itemize}
\item scarso interesse da parte dei membri del gruppo di lavoro nella tecnologia TTS;
\item difficoltà nel trovare un idea per un applicazione che fosse innovativa e utile;
\item le tecnologie necessarie alla realizzazione del capitolato non sono conosciute dai membri del gruppo e c'é il rischio di non riuscire ad avere una stima approssimativa della mole di lavoro da svolgere;
\item mancanza di conoscenze nell'implementazione di meccanismi atti a risolvere le problematiche legate all'utilizzo di un servizio remoto;
\item l'obbiettivo di rendere disponibile la sintesi vocale a tutte le applicazioni e fornire un interfaccia di configurazione estesa comporta conoscenze specifiche e uno studio particolarmente approfondito sulle modalità di interazione tra applicazioni e funzionalità del sistema operativo.
\end{itemize}

\end{document}
